\chapter{Missing data, structural zeros, and what is not implemented}\label{ch:missing}

The first part of this chapter explains how \nwsaom{} treats data that are unknown (missing ties and behavior values) or fixed by design (structural zeros and ones). The second part is an honest list of what RSiena offers and \nwsaom{} does not.

\section{Missing data: \opt{missnet()} and \opt{missbeh()}}\label{sec:missing}
Missing data here means dyads or actors whose true value at a wave is \emph{unknown}. This is RSiena's regular missing-data mechanism and it is different from composition change (actors who are not part of the network at a wave, \opt{present()} or \opt{joinat()}/\opt{leaveat()}): a missing tie is uncertainty about an otherwise active actor. The two can be used together.

\paragraph{Specifying missingness.}
\begin{itemize}
\item \opt{missnet(matlist)} takes one 0/1 $n\times n$ \emph{matrix} name per wave, in the order of \opt{wave1()}/\opt{wave2()} or \opt{waves()}; a 1 marks a dyad missing at that wave. It is a raw Stata matrix, not an \texttt{nwset} network. Build it with \texttt{matrix input} or \texttt{mkmat} from your own indicator. With two waves, \texttt{missnet(m1 m2)}.
\item \opt{missbeh(varlist)} takes one 0/1 \emph{variable} per wave (the same one-variable-per-wave convention as \opt{present()}), marking the actors whose behavior value is missing at that wave. It requires \opt{behavior()}.
\item Both are optional and independent: use either or both.
\end{itemize}
\begin{lstlisting}
nwwebuse glasgow, nwclear
* the first five actors did not report their ties at wave 2
matrix m1 = J(50,50,0)
matrix m2 = J(50,50,0)
forvalues i = 1/5 {
    forvalues j = 1/50 {
        if `i' != `j' matrix m2[`i',`j'] = 1
    }
}
nwsaom, wave1(glasgow1) wave2(glasgow2) outdegree reciprocity transtrip ///
    missnet(m1 m2) seed(1)
\end{lstlisting}
(This example was run with the development build; it converged, with an overall convergence ratio of 0.119.) Because the missing indicators are Stata matrices, they are limited by Stata's matrix size: 800$\times$800 in Stata/BE, 11,000$\times$11,000 in Stata/SE.

\paragraph{How missing values are handled.} \nwsaom{} follows RSiena's two steps.
\begin{enumerate}
\item \textbf{Imputation.} Each missing dyad is filled in by carrying forward the last observed value (0 if never observed up to that wave). A missing behavior value is filled in from the previous observation, else the next observation, else the mode of that wave (cross-sectional). The imputed values take part in the simulation completely normally; unlike \opt{present()}, missingness does not restrict who can act.
\item \textbf{Masking.} A dyad or actor missing at \emph{either} end wave of a period is excluded from the observed target statistic \emph{and} from every simulated statistic of that period, so the moment condition is not biased by the exclusion. The same applies to the distance that conditional estimation simulates to, and to the distance statistic of the rate under unconditional estimation. The existing effects needed no special-casing: their statistic functions are reused unchanged.
\end{enumerate}
Missing data are handled in the native engine, when the model is otherwise native-eligible; the masked statistic matches the Mata engine to machine precision.

\paragraph{A disclosed approximation.} The statistic of \opt{avalt} and \opt{avsim} for an actor depends on the \emph{current} values of the actor's alters, which differ across simulated replicates. For these effects the masking is only an approximation (it is exact for \opt{linear} and \opt{quadratic}). It is empirically better than not masking, but may be less precise than an idealised per-effect masking under heavy missingness combined with \opt{avalt} or \opt{avsim}.

\paragraph{Missing covariate values.} Covariates of \opt{nodematch()}, \opt{nodecov()}, \opt{nodeicov()}, \opt{nodeocov()}, \opt{simcov()} and their aliases follow RSiena. A missing value is imputed with the covariate mean when actors decide their ministeps, but the \emph{statistic} (target and every simulated statistic) skips a tie whose relevant actor is missing: the alter for \opt{altx()}, the ego for \opt{egox()}, either for \opt{samex()} and \opt{simx()}. With the default centring a missing \opt{altx()}, \opt{egox()} or \opt{nodecov()} value becomes 0 and skipping changes nothing; under \opt{nocenter} and for \opt{samex()}/\opt{simx()} it does, and the range and similarity mean of \opt{simx()} use the observed values only. On RSiena's \texttt{s501}/\texttt{s502} with 8 of 50 covariate values missing, density, reciprocity, altX, egoX and simX agree with RSiena 1.6.6 within sampling noise (simX $0.79$ against $0.78$). An \opt{interact()} naming a \opt{samex()} or \opt{simx()} effect whose variable has missing values is rejected.

\section{Structural zeros and ones: \opt{structural()}}\label{sec:structural}
A structural dyad is one whose tie value is known and cannot change, for example a legally mandated reporting tie, a physically impossible tie, or a dyad that you wish to hold fixed for a counterfactual. It is a simpler idea than missing data: the value is not unknown. RSiena marks such dyads by embedding the sentinel values 10 (structural zero) and 11 (structural one) in the network array. \nwsaom{} uses a \emph{separate} 0/1 mask matrix instead, in the manner of \opt{missnet()}.
\begin{itemize}
\item \opt{structural(matname)} takes \emph{one} 0/1 $n\times n$ matrix with a zero diagonal; a 1 freezes that dyad for the whole period.
\item The observed value of a marked dyad must be \emph{identical} at both waves. A dyad that genuinely changed cannot be structural, and \nwsaom{} rejects it with an error (rc 198) rather than ignore it.
\item Mechanically the dyad is excluded from every actor's candidate set in each ministep, so it can never toggle; the effect statistics and their changes are reused unchanged.
\end{itemize}
The following example marks dyads that are zero in both waves among the first five actors:
\begin{lstlisting}
nwtomata glasgow1, mat(A1)
nwtomata glasgow2, mat(A2)
mata:
S = J(50,50,0)
S[|1,1\5,50|] = (A1[|1,1\5,50|]:==0):*(A2[|1,1\5,50|]:==0)
S = S - diag(diagonal(S))
st_matrix("S", S)
end
nwsaom, wave1(glasgow1) wave2(glasgow2) outdegree reciprocity transtrip ///
    structural(S) seed(1)
\end{lstlisting}
(Also run with the development build; the fit converged with an overall ratio of 0.110.) The structural mask also underlies the two-mode implementation, where only actor--event dyads are free.

\paragraph{Scope.} \opt{structural()} is implemented for the two-wave path \opt{wave1()}/\opt{wave2()}, for network-only models, and it is not combinable with \opt{symmetric}, \opt{ratecov()} or the network endowment/creation split; each of these combinations is rejected outright. For users the option is limited to the two-wave path; it is used with \opt{waves()} only internally, for two-mode networks. It runs in the native plugin.

\section{What RSiena has that \nwsaom{} does not (yet)}\label{sec:notimpl}
RSiena's effects table, as returned by \texttt{getEffects()} for 12 data configurations of RSiena 1.6.14, has \nTotal{} rows (a row is one RSiena short name for one variable type and interaction context). Every row was compared with \nwsaom{} by reading the code and the documentation, and by running models against \texttt{siena07} where they were implemented. Of these rows:
\begin{itemize}
\item \nImpl{} are implemented, \nPartial{} is partly implemented;
\item \nNotImpl{} are \emph{not implemented};
\item \nNotPlanned{} are deliberately out of scope: continuous (SDE) behavior and the effects that involve a continuous dependent variable.
\end{itemize}
Of the \nPlannable{} rows that are not out of scope, \pctImpl\% are implemented. The percentage is by \emph{rows}, not by how often an effect is used: the effects most often used in applied papers are largely covered, and the missing rows are concentrated in a few groups. The groups follow; the row-level list is at the end of the chapter.

\subsection{Effects not implemented, by reason}
The \nImpl{} implemented rows are \pctImpl\% of the plannable ones. The \nNotImpl{} rows that are not implemented are of five kinds.
\begin{description}[leftmargin=0pt,labelindent=0pt]
\item[The generalised method of moments (31 rows).] RSiena's GMoM estimator uses its own statistics (the effects named like \texttt{egoX\_gmm}, \texttt{avAlt\_gmm}, \texttt{crprod\_gmm}, and related placeholders). \nwsaom{} implements the standard method of moments, and maximum likelihood and Bayesian estimation, but not the generalised method of moments.
\item[The settings model (12 rows).] The \texttt{setting*} effects (\texttt{settingSizeAct}, \texttt{settingOppAct} and relatives) belong to RSiena's settings model, a separate model family with setting steps and a rate for each setting; \nwsaom{} has the standard tie-change model only. There is also nothing to validate against by estimation: in \texttt{siena07} eight of the twelve effects have a statistic that is constant under simulation (the estimate stays 0 and has no standard error), and the four variants with the suffix D have a zero derivative and are fixed at 0.
\item[Effects that RSiena cannot evaluate or does not offer (15 rows).] \texttt{siena07} stops or returns nothing useful for these combinations, so there is nothing to compare with and \nwsaom{} refuses them with an error message (Table below and Section~\ref{sec:cannoteval}).
\item[Placeholders (2 rows).] \texttt{minXAlt} and \texttt{maxXAlt} in a non-directed network with unspecified variables are entries of RSiena's effects table that \texttt{siena07} does not offer.
\item[Two further effects.] \texttt{sameWXClosure} for a two-mode network together with a co-evolving behavior inside \texttt{nwsaom multinet} (a two-mode member and a behavior cannot be combined there), and \texttt{RateX} for a continuous behavior variable (continuous behavior is not planned).
\end{description}
\noindent In addition \nNotPlanned{} rows are out of scope: continuous (stochastic differential equation) dependent variables and all effects that involve them.

\subsection{Effects that RSiena cannot evaluate}\label{sec:cannoteval}
\begin{center}\small
\begin{tabular}{@{}p{4.2cm}p{10.4cm}@{}}
\toprule
Effect & What \texttt{siena07} does\\
\midrule
\texttt{inPop\_ego}, two-mode network & stops with a C++ exception\\
\texttt{sameXOutAct}, \texttt{diffXOutAct}, two-mode network, with a constant covariate, a changing covariate or a behavior & the R process aborts\\
\texttt{avGroupEgoX}, \texttt{totGroupEgoX}, two-mode network, constant covariate & reports the effect as not meaningful (available with a changing covariate)\\
\texttt{homXOutAct2} with a behavior & reports the effect as not implemented for behavioral variables\\
\texttt{inRateInv}, \texttt{inRateLog}, two-mode network & stops with an error (the receivers of ties are not the actors)\\
\texttt{susceptAvCovar}, \texttt{infectCovar} with a behavior as covariate & no such effect exists (the covariate is a constant or changing covariate)\\
\texttt{outOutDist2AvIntn} & listed in the effects table, but there is no implementation\\
\texttt{toU} for a one-mode network that depends on a two-mode one & rejected\\
\texttt{avAlt.2M.tie} with RSiena's default parameter 3 & throws an error; \nwsaom{} offers the parameters 1 and 2\\
\bottomrule
\end{tabular}
\end{center}
Further, \opt{higher} in a non-directed network equals the density statistic and is not identified: RSiena returns a standard error of NA and \nwsaom{} stops.

\subsection{RSiena features beyond effects}\label{sec:est-notimpl}
Besides effects, RSiena provides estimation and analysis tools. This is how they compare.
\begin{description}[leftmargin=0pt,labelindent=0pt]
\item[Method of moments (\texttt{siena07}).] Implemented, conditional and unconditional, with the same phase structure (Chapter~\ref{ch:estimation}).
\item[Maximum likelihood and Bayesian estimation (\texttt{siena07} with \texttt{maxlike = TRUE}, \texttt{multi\_siena}).] Available for a directed one-mode network, with or without one co-evolving behavior: the option \opt{ml} and the subcommand \texttt{nwsaom bayes} (Chapter~\ref{ch:estimation}). Not available for non-directed and two-mode networks, several networks, missing data, composition change and rate covariates. The hierarchical Bayesian analysis of several groups (\texttt{multi\_siena} with several groups) and Bayesian estimation of a behavior alone are not implemented.
\item[Multi-group analysis (\texttt{siena08}).] The meta-analysis of separate fits is available as \texttt{nwsaom meta} (random-effects pooling by robust IWLS, heterogeneity tests, Fisher combination, maximum-likelihood mean and between-group sd with profile intervals). RSiena's joint multi-group \emph{estimation} (\texttt{sienaGroupCreate}, one model fitted to several groups at once) is not offered: fit each group separately and pool.
\item[Goodness of fit (\texttt{sienaGOF}).] Available as \texttt{estat gof} (Chapter~\ref{ch:estimation}), with the differences listed there, and not for fits with several co-evolving behaviors.
\item[Composition change.] Available, including fractional join/leave times (\opt{joinat()}/\opt{leaveat()}) as in \texttt{sienaCompositionChange()}. Composition change and missing ties are not available inside \texttt{nwsaom multinet}.
\item[Continuous behavior.] Not implemented (out of scope).
\item[Simulation only, time-heterogeneity tests, score-type tests, effect importance.] Available: \texttt{nwsaom simulate}, \texttt{estat timetest}, \texttt{estat score} and \texttt{estat importance} (Chapter~\ref{ch:tools}). After \texttt{nwsaom multinet} the tests are available, but not \texttt{estat gof}, \texttt{estat importance} and \texttt{estat simulate}.
\item[Fixing and testing individual parameters.] Available: \opt{fix()} holds an effect at its \opt{theta0()} value (RSiena's \texttt{fix}), and \texttt{estat score} tests it, also after \texttt{nwsaom multinet}.
\item[The generalised method of moments.] Not implemented (31 statistics).
\item[The settings model.] Not implemented (12 effects).
\item[Model types for non-directed relations.] Offered by name or RSiena's modelType number (\opt{symtype()}: forcing, confirmation, force, agree, joint; modelType 2, 3, 4, 5, 6).
\end{description}

\section{Row-level list of the effects that are not implemented}\label{sec:rowlist}
The tables below list every RSiena effects-table row that is \emph{not implemented} (one table per reason), in the order of RSiena's effects table; the final entry is the out-of-scope group. A row names an RSiena short name together with the variable type and the interaction context (another network, a constant or changing covariate, a dyadic covariate, a behavior variable). Rows that are implemented are listed in the appendix on effects.

\subsection*{Cross-network effects that also involve a behavior variable (needs behavior inside nwsaom multinet) (37 rows)}
\footnotesize
\begin{longtable}{@{}p{2.8cm}p{4.6cm}p{6.5cm}@{}}
\toprule
RSiena name & Network / context & Note\\
\midrule
\endhead
\texttt{sameWWClosure} & one-mode network + another network + a behavior variable & Not implemented.\\
\texttt{diffWWClosure} & one-mode network + another network + a behavior variable & Not implemented.\\
\texttt{jumpWWClosure} & one-mode network + another network + a behavior variable & Not implemented.\\
\texttt{jumpFrom} & one-mode network + another network + a behavior variable & Not implemented.\\
\texttt{jumpSharedIn} & one-mode network + another network + a behavior variable & Not implemented.\\
\texttt{sameWXClosure} & one-mode network + another network + a behavior variable & Not implemented.\\
\texttt{diffWXClosure} & one-mode network + another network + a behavior variable & Not implemented.\\
\texttt{jumpWXClosure} & one-mode network + another network + a behavior variable & Not implemented.\\
\texttt{homWXClosure} & one-mode network + another network + a behavior variable & Not implemented.\\
\texttt{sameXWClosure} & one-mode network + another network + a behavior variable & Not implemented.\\
\texttt{diffXWClosure} & one-mode network + another network + a behavior variable & Not implemented.\\
\texttt{homCovNetNet} & one-mode network + another network + a behavior variable & Not implemented.\\
\texttt{contrastCovNetNet} & one-mode network + another network + a behavior variable & Not implemented.\\
\texttt{allDifCovNetNet} & one-mode network + another network + a behavior variable & Not implemented.\\
\texttt{homCovNetNetIn} & one-mode network + another network + a behavior variable & Not implemented.\\
\texttt{contrastCovNetNetIn} & one-mode network + another network + a behavior variable & Not implemented.\\
\texttt{allDifCovNetNetIn} & one-mode network + another network + a behavior variable & Not implemented.\\
\texttt{simEgoDist2W} & one-mode network + another network + a behavior variable & Not implemented.\\
\texttt{altInDist2W} & one-mode network + another network + a behavior variable & Not implemented.\\
\texttt{totInDist2W} & one-mode network + another network + a behavior variable & Not implemented.\\
\texttt{inActIntnX} & one-mode network + another network + a behavior variable & Not implemented.\\
\texttt{sameXInActIntn} & one-mode network + another network + a behavior variable & Not implemented.\\
\texttt{covNetNet} & one-mode network + another network + a behavior variable & Not implemented.\\
\texttt{covNetNetIn} & one-mode network + another network + a behavior variable & Not implemented.\\
\texttt{inPopIntnX} & one-mode network + another network + a behavior variable & Not implemented.\\
\texttt{sameXInPopIntn} & one-mode network + another network + a behavior variable & Not implemented.\\
\texttt{simEgoInDist2W} & one-mode network + another network + a behavior variable & Not implemented.\\
\texttt{outPopIntnX} & one-mode network + another network + a behavior variable & Not implemented.\\
\texttt{outActIntnX} & one-mode network + another network + a behavior variable & Not implemented.\\
\texttt{sameXOutPopIntn} & one-mode network + another network + a behavior variable & Not implemented.\\
\texttt{sameXOutActIntn} & one-mode network + another network + a behavior variable & Not implemented.\\
\texttt{altDist2W} & one-mode network + another network + a behavior variable & Not implemented.\\
\texttt{totDist2W} & one-mode network + another network + a behavior variable & Not implemented.\\
\texttt{simDist2W} & one-mode network + another network + a behavior variable & Not implemented.\\
\texttt{sameWXClosure} & two-mode network + another network + a behavior variable & Not implemented.\\
\texttt{susceptAvCovar} & rate (behavior) + another network + a behavior variable & Structural/exposure/diffusion rate effects are not implemented; only a constant rate and ratecov() exist.\\
\texttt{infectCovar} & rate (behavior) + another network + a behavior variable & Structural/exposure/diffusion rate effects are not implemented; only a constant rate and ratecov() exist.\\
\bottomrule
\end{longtable}
\normalsize

\subsection*{Rate effects (covariates on the rate of two-mode networks and behaviors, exposure and infection effects) (36 rows)}
\footnotesize
\begin{longtable}{@{}p{2.8cm}p{4.6cm}p{6.5cm}@{}}
\toprule
RSiena name & Network / context & Note\\
\midrule
\endhead
\texttt{RateX} & rate (one-mode network) + changing covariate & ratecov() takes constant covariates; tvcov() names are not resolved for ratecov.\\
\texttt{RateX} & rate (one-mode network) + a behavior variable & ratecov() is a network-rate option for constant node covariates only.\\
\texttt{outRate} & rate (behavior) + another network & Structural/exposure/diffusion rate effects are not implemented; only a constant rate and ratecov() exist.\\
\texttt{inRate} & rate (behavior) + another network & Structural/exposure/diffusion rate effects are not implemented; only a constant rate and ratecov() exist.\\
\texttt{recipRate} & rate (behavior) + another network & Structural/exposure/diffusion rate effects are not implemented; only a constant rate and ratecov() exist.\\
\texttt{avExposure} & rate (behavior) + another network & Structural/exposure/diffusion rate effects are not implemented; only a constant rate and ratecov() exist.\\
\texttt{susceptAvIn} & rate (behavior) + another network & Structural/exposure/diffusion rate effects are not implemented; only a constant rate and ratecov() exist.\\
\texttt{totExposure} & rate (behavior) + another network & Structural/exposure/diffusion rate effects are not implemented; only a constant rate and ratecov() exist.\\
\texttt{infectIn} & rate (behavior) + another network & Structural/exposure/diffusion rate effects are not implemented; only a constant rate and ratecov() exist.\\
\texttt{infectOut} & rate (behavior) + another network & Structural/exposure/diffusion rate effects are not implemented; only a constant rate and ratecov() exist.\\
\texttt{avTinExposureDist2} & rate (behavior) + another network & Structural/exposure/diffusion rate effects are not implemented; only a constant rate and ratecov() exist.\\
\texttt{totInExposureDist2} & rate (behavior) + another network & Structural/exposure/diffusion rate effects are not implemented; only a constant rate and ratecov() exist.\\
\texttt{totAInExposureDist2} & rate (behavior) + another network & Structural/exposure/diffusion rate effects are not implemented; only a constant rate and ratecov() exist.\\
\texttt{anyInExposureDist2} & rate (behavior) + another network & Structural/exposure/diffusion rate effects are not implemented; only a constant rate and ratecov() exist.\\
\texttt{RateX} & rate (behavior) + constant covariate & ratecov() is a network-rate option for constant node covariates only.\\
\texttt{susceptAvCovar} & rate (behavior) + another network + constant covariate & Structural/exposure/diffusion rate effects are not implemented; only a constant rate and ratecov() exist.\\
\texttt{infectCovar} & rate (behavior) + another network + constant covariate & Structural/exposure/diffusion rate effects are not implemented; only a constant rate and ratecov() exist.\\
\texttt{RateX} & rate (behavior) + a behavior variable & ratecov() is a network-rate option for constant node covariates only.\\
\texttt{RateX} & rate (behavior) + changing covariate & ratecov() takes constant covariates; tvcov() names are not resolved for ratecov.\\
\texttt{susceptAvCovar} & rate (behavior) + another network + changing covariate & Structural/exposure/diffusion rate effects are not implemented; only a constant rate and ratecov() exist.\\
\texttt{infectCovar} & rate (behavior) + another network + changing covariate & Structural/exposure/diffusion rate effects are not implemented; only a constant rate and ratecov() exist.\\
\texttt{outRate} & rate (two-mode network) & Structural/exposure/diffusion rate effects are not implemented; only a constant rate and ratecov() exist.\\
\texttt{outRateInv} & rate (two-mode network) & Structural/exposure/diffusion rate effects are not implemented; only a constant rate and ratecov() exist.\\
\texttt{outRateLog} & rate (two-mode network) & Structural/exposure/diffusion rate effects are not implemented; only a constant rate and ratecov() exist.\\
\texttt{inRateInv} & rate (two-mode network) & Structural/exposure/diffusion rate effects are not implemented; only a constant rate and ratecov() exist.\\
\texttt{inRateLog} & rate (two-mode network) & Structural/exposure/diffusion rate effects are not implemented; only a constant rate and ratecov() exist.\\
\texttt{RateX} & rate (two-mode network) + constant covariate & ratecov() is a network-rate option for constant node covariates only.\\
\texttt{RateX} & rate (two-mode network) + changing covariate & ratecov() takes constant covariates; tvcov() names are not resolved for ratecov.\\
\texttt{RateX} & rate (two-mode network) + a behavior variable & ratecov() is a network-rate option for constant node covariates only.\\
\texttt{RateX} & rate (one-mode network) + continuous variable & ratecov() is a network-rate option for constant node covariates only.\\
\texttt{outRate} & rate (non-directed network) & Structural/exposure/diffusion rate effects are not implemented; only a constant rate and ratecov() exist.\\
\texttt{outRateInv} & rate (non-directed network) & Structural/exposure/diffusion rate effects are not implemented; only a constant rate and ratecov() exist.\\
\texttt{outRateLog} & rate (non-directed network) & Structural/exposure/diffusion rate effects are not implemented; only a constant rate and ratecov() exist.\\
\texttt{RateX} & rate (non-directed network) + changing covariate & ratecov() takes constant covariates; tvcov() names are not resolved for ratecov. Validated 2026-10-23: plain 3-wave pairwise fit agrees with siena07 to three d…\\
\texttt{RateX} & rate (non-directed network) + a behavior variable & ratecov() is a network-rate option for constant node covariates only. Validated 2026-10-23: plain 3-wave pairwise fit agrees with siena07 to three digits\\
\texttt{infectDeg} & rate (behavior (non-directed network)) + another network & Structural/exposure/diffusion rate effects are not implemented; only a constant rate and ratecov() exist.\\
\bottomrule
\end{longtable}
\normalsize

\subsection*{Other (see the effect notes) (32 rows)}
\footnotesize
\begin{longtable}{@{}p{2.8cm}p{4.6cm}p{6.5cm}@{}}
\toprule
RSiena name & Network / context & Note\\
\midrule
\endhead
\texttt{inPop\_ego} & two-mode network & Not implemented.\\
\texttt{unspInt} & two-mode network & Two-mode: interact() is refused.\\
\texttt{unspInt} & non-directed network & interact() is refused with symmetric (tvcov/interact guard).\\
\texttt{sameXOutAct} & two-mode network + constant covariate & Other covariate effects are refused for two-mode (refused with an error message).\\
\texttt{diffXOutAct} & two-mode network + constant covariate & Other covariate effects are refused for two-mode (refused with an error message).\\
\texttt{avGroupEgoX} & two-mode network + constant covariate & Other covariate effects are refused for two-mode (refused with an error message).\\
\texttt{totGroupEgoX} & two-mode network + constant covariate & Other covariate effects are refused for two-mode (refused with an error message).\\
\texttt{sameXOutAct} & two-mode network + changing covariate & tvcov() is refused with two-mode networks.\\
\texttt{diffXOutAct} & two-mode network + changing covariate & tvcov() is refused with two-mode networks.\\
\texttt{higher} & non-directed network + changing covariate & tvcov() is refused with symmetric. Validated 2026-10-23: observed statistic equals siena07 and estimates agree over 2 seeds (target only; the effect is refused…\\
\texttt{X} & two-mode network + constant dyadic covariate & dyadcov() is refused for two-mode.\\
\texttt{XRecip} & non-directed network + changing dyadic covariate & dyadcov() is time-constant only (a single matrix); no changing dyadic covariate. Validated 2026-10-23: observed statistic equals siena07 and estimates agree ov…\\
\texttt{homXOutAct2} & one-mode network + a behavior variable & Only simX of the dependent behavior is available (behsim); nodeocov()/nodeicov() take fixed Stata variables, not the evolving behavior.\\
\texttt{sameXOutAct} & two-mode network + a behavior variable & Only simX of the dependent behavior is available (behsim); nodeocov()/nodeicov() take fixed Stata variables, not the evolving behavior.\\
\texttt{diffXOutAct} & two-mode network + a behavior variable & Only simX of the dependent behavior is available (behsim); nodeocov()/nodeicov() take fixed Stata variables, not the evolving behavior.\\
\texttt{crossXOutAct} & two-mode network + a behavior variable & Only simX of the dependent behavior is available (behsim); nodeocov()/nodeicov() take fixed Stata variables, not the evolving behavior.\\
\texttt{avGroupEgoX} & two-mode network + a behavior variable & Only simX of the dependent behavior is available (behsim); nodeocov()/nodeicov() take fixed Stata variables, not the evolving behavior.\\
\texttt{totGroupEgoX} & two-mode network + a behavior variable & Only simX of the dependent behavior is available (behsim); nodeocov()/nodeicov() take fixed Stata variables, not the evolving behavior.\\
\texttt{degAbsDiffX} & non-directed network + a behavior variable & Only simX of the dependent behavior is available (behsim); nodeocov()/nodeicov() take fixed Stata variables, not the evolving behavior.\\
\texttt{degPosDiffX} & non-directed network + a behavior variable & Only simX of the dependent behavior is available (behsim); nodeocov()/nodeicov() take fixed Stata variables, not the evolving behavior.\\
\texttt{higher} & non-directed network + a behavior variable & Only simX of the dependent behavior is available (behsim); nodeocov()/nodeicov() take fixed Stata variables, not the evolving behavior. Validated 2026-10-23: o…\\
\texttt{homXOutAct2} & non-directed network + a behavior variable & Only simX of the dependent behavior is available (behsim); nodeocov()/nodeicov() take fixed Stata variables, not the evolving behavior.\\
\texttt{from.w.ind} & one-mode network + another network + another network & Not implemented.\\
\texttt{betweenPop} & one-mode network + another network & Cross-network effect: nwsaom multiplex supports only outdegree, reciprocity and crprod (two directed networks, two waves).\\
\texttt{toAny} & one-mode network + another network & Cross-network effect: nwsaom multiplex supports only outdegree, reciprocity and crprod (two directed networks, two waves).\\
\texttt{toU} & one-mode network + another network + constant dyadic covariate & Not implemented.\\
\texttt{constant} & behavior & nwsaom behavior(): only linear, quadratic, avalt, avsim (plus endow/creation of quadratic/avalt/avsim) are available.\\
\texttt{behUnspInt} & behavior & interact() does not support behavior interactions.\\
\texttt{avGroupEgoX} & behavior + a behavior variable & nwsaom has no covariate or cross-behavior effects on a behavior (several behaviors interact only through the network).\\
\texttt{totGroupEgoX} & behavior + a behavior variable & nwsaom has no covariate or cross-behavior effects on a behavior (several behaviors interact only through the network).\\
\texttt{minXAlt} & non-directed network + ?yyyyyy + ?zzzzzz & Placeholder row: RSiena lists this effect for a non-directed network with unresolved interaction placeholders (interaction1 = yyyyyy); not a usable effect with…\\
\texttt{maxXAlt} & non-directed network + ?yyyyyy + ?zzzzzz & Placeholder row: RSiena lists this effect for a non-directed network with unresolved interaction placeholders (interaction1 = yyyyyy); not a usable effect with…\\
\bottomrule
\end{longtable}
\normalsize

\subsection*{GMoM statistics (the generalised method of moments estimator is not implemented) (31 rows)}
\footnotesize
\begin{longtable}{@{}p{2.8cm}p{4.6cm}p{6.5cm}@{}}
\toprule
RSiena name & Network / context & Note\\
\midrule
\endhead
\texttt{newrecip} & gmm (one-mode network) & GMoM statistic: nwsaom estimates by Method of Moments only (no GMoM estimator).\\
\texttt{persistrecip} & gmm (one-mode network) & GMoM statistic: nwsaom estimates by Method of Moments only (no GMoM estimator).\\
\texttt{realrecip} & gmm (one-mode network) & GMoM statistic: nwsaom estimates by Method of Moments only (no GMoM estimator).\\
\texttt{realtrans} & gmm (one-mode network) & GMoM statistic: nwsaom estimates by Method of Moments only (no GMoM estimator).\\
\texttt{agreetrans} & gmm (one-mode network) & GMoM statistic: nwsaom estimates by Method of Moments only (no GMoM estimator).\\
\texttt{egoX\_gmm} & gmm (one-mode network) + constant covariate & GMoM statistic: nwsaom estimates by Method of Moments only (no GMoM estimator).\\
\texttt{simX\_gmm} & gmm (one-mode network) + constant covariate & GMoM statistic: nwsaom estimates by Method of Moments only (no GMoM estimator).\\
\texttt{egoX\_gmm} & gmm (one-mode network) + changing covariate & GMoM statistic: nwsaom estimates by Method of Moments only (no GMoM estimator).\\
\texttt{simX\_gmm} & gmm (one-mode network) + changing covariate & GMoM statistic: nwsaom estimates by Method of Moments only (no GMoM estimator).\\
\texttt{egoX\_gmm} & gmm (one-mode network) + a behavior variable & GMoM statistic: nwsaom estimates by Method of Moments only (no GMoM estimator).\\
\texttt{simX\_gmm} & gmm (one-mode network) + a behavior variable & GMoM statistic: nwsaom estimates by Method of Moments only (no GMoM estimator).\\
\texttt{avSim\_gmm} & gmm (behavior) + another network & GMoM statistic: nwsaom estimates by Method of Moments only (no GMoM estimator).\\
\texttt{avAlt\_gmm} & gmm (behavior) + another network & GMoM statistic: nwsaom estimates by Method of Moments only (no GMoM estimator).\\
\texttt{totAlt\_gmm} & gmm (behavior) + another network & GMoM statistic: nwsaom estimates by Method of Moments only (no GMoM estimator).\\
\texttt{maxAlt\_gmm} & gmm (behavior) + another network & GMoM statistic: nwsaom estimates by Method of Moments only (no GMoM estimator).\\
\texttt{minAlt\_gmm} & gmm (behavior) + another network & GMoM statistic: nwsaom estimates by Method of Moments only (no GMoM estimator).\\
\texttt{crprod\_gmm} & gmm (one-mode network) + another network & GMoM statistic: nwsaom estimates by Method of Moments only (no GMoM estimator).\\
\texttt{from\_gmm} & gmm (one-mode network) + another network & GMoM statistic: nwsaom estimates by Method of Moments only (no GMoM estimator).\\
\texttt{to\_gmm} & gmm (one-mode network) + another network & GMoM statistic: nwsaom estimates by Method of Moments only (no GMoM estimator).\\
\texttt{to\_gmm} & gmm (two-mode network) + another network & GMoM statistic: nwsaom estimates by Method of Moments only (no GMoM estimator).\\
\texttt{egoX\_gmm} & gmm (one-mode network) + continuous variable & GMoM statistic: nwsaom estimates by Method of Moments only (no GMoM estimator).\\
\texttt{simX\_gmm} & gmm (one-mode network) + continuous variable & GMoM statistic: nwsaom estimates by Method of Moments only (no GMoM estimator).\\
\texttt{egoX\_gmm} & gmm (non-directed network) + constant covariate & GMoM statistic: nwsaom estimates by Method of Moments only (no GMoM estimator).\\
\texttt{simX\_gmm} & gmm (non-directed network) + constant covariate & GMoM statistic: nwsaom estimates by Method of Moments only (no GMoM estimator).\\
\texttt{egoX\_gmm} & gmm (non-directed network) + changing covariate & GMoM statistic: nwsaom estimates by Method of Moments only (no GMoM estimator).\\
\texttt{simX\_gmm} & gmm (non-directed network) + changing covariate & GMoM statistic: nwsaom estimates by Method of Moments only (no GMoM estimator).\\
\texttt{egoX\_gmm} & gmm (non-directed network) + a behavior variable & GMoM statistic: nwsaom estimates by Method of Moments only (no GMoM estimator).\\
\texttt{simX\_gmm} & gmm (non-directed network) + a behavior variable & GMoM statistic: nwsaom estimates by Method of Moments only (no GMoM estimator).\\
\texttt{from\_gmm} & gmm (non-directed network) + another network & GMoM statistic: nwsaom estimates by Method of Moments only (no GMoM estimator).\\
\texttt{to\_gmm} & gmm (non-directed network) + another network & GMoM statistic: nwsaom estimates by Method of Moments only (no GMoM estimator).\\
\texttt{crprod\_gmm} & gmm (non-directed network) + another network & GMoM statistic: nwsaom estimates by Method of Moments only (no GMoM estimator).\\
\bottomrule
\end{longtable}
\normalsize

\subsection*{Several networks with a two-mode member (nwsaom multinet supports directed and non-directed members only) (25 rows)}
\footnotesize
\begin{longtable}{@{}p{2.8cm}p{4.6cm}p{6.5cm}@{}}
\toprule
RSiena name & Network / context & Note\\
\midrule
\endhead
\texttt{inActIntn} & two-mode network + another network & Cross-network effect: nwsaom multiplex supports only outdegree, reciprocity and crprod (two directed networks, two waves).\\
\texttt{outActIntn} & two-mode network + another network & Cross-network effect: nwsaom multiplex supports only outdegree, reciprocity and crprod (two directed networks, two waves).\\
\texttt{divOutEgoIntn} & two-mode network + another network & Cross-network effect: nwsaom multiplex supports only outdegree, reciprocity and crprod (two directed networks, two waves).\\
\texttt{divInEgoIntn} & two-mode network + another network & Cross-network effect: nwsaom multiplex supports only outdegree, reciprocity and crprod (two directed networks, two waves).\\
\texttt{divInAltIntn} & two-mode network + another network & Cross-network effect: nwsaom multiplex supports only outdegree, reciprocity and crprod (two directed networks, two waves).\\
\texttt{avDegIntn} & two-mode network + another network & Cross-network effect: nwsaom multiplex supports only outdegree, reciprocity and crprod (two directed networks, two waves).\\
\texttt{outOutActIntn} & two-mode network + another network & Cross-network effect: nwsaom multiplex supports only outdegree, reciprocity and crprod (two directed networks, two waves).\\
\texttt{outOutAvIntn} & two-mode network + another network & Cross-network effect: nwsaom multiplex supports only outdegree, reciprocity and crprod (two directed networks, two waves).\\
\texttt{dist2OutInActIntn} & two-mode network + another network & Cross-network effect: nwsaom multiplex supports only outdegree, reciprocity and crprod (two directed networks, two waves).\\
\texttt{outOutDist2ActIntn} & two-mode network + another network & Cross-network effect: nwsaom multiplex supports only outdegree, reciprocity and crprod (two directed networks, two waves).\\
\texttt{avAlt.2M.tot} & two-mode network + another network & Cross-network effect: nwsaom multiplex supports only outdegree, reciprocity and crprod (two directed networks, two waves).\\
\texttt{to} & two-mode network + another network & Cross-network effect: nwsaom multiplex supports only outdegree, reciprocity and crprod (two directed networks, two waves).\\
\texttt{toBack} & two-mode network + another network & Cross-network effect: nwsaom multiplex supports only outdegree, reciprocity and crprod (two directed networks, two waves).\\
\texttt{toRecip} & two-mode network + another network & Cross-network effect: nwsaom multiplex supports only outdegree, reciprocity and crprod (two directed networks, two waves).\\
\texttt{sharedTo} & two-mode network + another network & Cross-network effect: nwsaom multiplex supports only outdegree, reciprocity and crprod (two directed networks, two waves).\\
\texttt{avAlt.2M.tie} & two-mode network + another network & Cross-network effect: nwsaom multiplex supports only outdegree, reciprocity and crprod (two directed networks, two waves).\\
\texttt{inPopOutW} & two-mode network + another network & Cross-network effect: nwsaom multiplex supports only outdegree, reciprocity and crprod (two directed networks, two waves).\\
\texttt{nDist2ActIntn} & two-mode network + another network & Cross-network effect: nwsaom multiplex supports only outdegree, reciprocity and crprod (two directed networks, two waves).\\
\texttt{toAny} & two-mode network + another network & Cross-network effect: nwsaom multiplex supports only outdegree, reciprocity and crprod (two directed networks, two waves).\\
\texttt{JoutMix} & two-mode network + another network & Cross-network effect: nwsaom multiplex supports only outdegree, reciprocity and crprod (two directed networks, two waves).\\
\texttt{gwespFBMix} & two-mode network + another network & Cross-network effect: nwsaom multiplex supports only outdegree, reciprocity and crprod (two directed networks, two waves).\\
\texttt{sameWXClosure} & two-mode network + another network + constant covariate & Not implemented.\\
\texttt{sameWXClosure} & two-mode network + another network + changing covariate & Not implemented.\\
\texttt{sharedToU} & two-mode network + another network + constant dyadic covariate & Not implemented.\\
\texttt{avAltU.2M.tie} & two-mode network + another network + constant dyadic covariate & Not implemented.\\
\bottomrule
\end{longtable}
\normalsize

\subsection*{RSiena's settings model (a different simulation procedure) (12 rows)}
\footnotesize
\begin{longtable}{@{}p{2.8cm}p{4.6cm}p{6.5cm}@{}}
\toprule
RSiena name & Network / context & Note\\
\midrule
\endhead
\texttt{settingSizeAct} & one-mode network & Not implemented.\\
\texttt{settingSizeActSqrt} & one-mode network & Not implemented.\\
\texttt{settingSizeActLog} & one-mode network & Not implemented.\\
\texttt{settingSizeActInv} & one-mode network & Not implemented.\\
\texttt{settingOppAct} & one-mode network & Not implemented.\\
\texttt{settingOppActSqrt} & one-mode network & Not implemented.\\
\texttt{settingOppActLog} & one-mode network & Not implemented.\\
\texttt{settingLogCreationAct} & one-mode network & Not implemented.\\
\texttt{settingOppActD} & one-mode network & Not implemented.\\
\texttt{settingOppActSqrtD} & one-mode network & Not implemented.\\
\texttt{settingOppActLogD} & one-mode network & Not implemented.\\
\texttt{settingLogCreationActD} & one-mode network & Not implemented.\\
\bottomrule
\end{longtable}
\normalsize

\subsection*{Non-directed network with another network: not validated or differs from siena07 (see the notes) (8 rows)}
\footnotesize
\begin{longtable}{@{}p{2.8cm}p{4.6cm}p{6.5cm}@{}}
\toprule
RSiena name & Network / context & Note\\
\midrule
\endhead
\texttt{from.w.ind} & non-directed network + another network + another network & Not implemented.\\
\texttt{toAny} & non-directed network + another network & Cross-network effect: nwsaom multiplex supports only outdegree, reciprocity and crprod (two directed networks, two waves).\\
\texttt{gwespMix} & non-directed network + another network & Cross-network effect: nwsaom multiplex supports only outdegree, reciprocity and crprod (two directed networks, two waves).\\
\texttt{outOutActIntn} & non-directed network + another network & Cross-network effect: nwsaom multiplex supports only outdegree, reciprocity and crprod (two directed networks, two waves).\\
\texttt{outOutAvIntn} & non-directed network + another network & Cross-network effect: nwsaom multiplex supports only outdegree, reciprocity and crprod (two directed networks, two waves).\\
\texttt{outOutDist2AvIntn} & non-directed network + another network & Cross-network effect: nwsaom multiplex supports only outdegree, reciprocity and crprod (two directed networks, two waves).\\
\texttt{both} & non-directed network + another network & Cross-network effect: nwsaom multiplex supports only outdegree, reciprocity and crprod (two directed networks, two waves).\\
\texttt{avAlt.2M.tie} & non-directed network + another network & Cross-network effect: nwsaom multiplex supports only outdegree, reciprocity and crprod (two directed networks, two waves).\\
\bottomrule
\end{longtable}
\normalsize

\subsection*{Behavior effects of two networks / further behavior-network effects (8 rows)}
\footnotesize
\begin{longtable}{@{}p{2.8cm}p{4.6cm}p{6.5cm}@{}}
\toprule
RSiena name & Network / context & Note\\
\midrule
\endhead
\texttt{FFDeg} & behavior + another network + another network & nwsaom behavior(): only linear, quadratic, avalt, avsim (plus endow/creation of quadratic/avalt/avsim) are available.\\
\texttt{BBDeg} & behavior + another network + another network & nwsaom behavior(): only linear, quadratic, avalt, avsim (plus endow/creation of quadratic/avalt/avsim) are available.\\
\texttt{FBDeg} & behavior + another network + another network & nwsaom behavior(): only linear, quadratic, avalt, avsim (plus endow/creation of quadratic/avalt/avsim) are available.\\
\texttt{FRDeg} & behavior + another network + another network & nwsaom behavior(): only linear, quadratic, avalt, avsim (plus endow/creation of quadratic/avalt/avsim) are available.\\
\texttt{BRDeg} & behavior + another network + another network & nwsaom behavior(): only linear, quadratic, avalt, avsim (plus endow/creation of quadratic/avalt/avsim) are available.\\
\texttt{indegMixedPop} & behavior + another network + another network & nwsaom behavior(): only linear, quadratic, avalt, avsim (plus endow/creation of quadratic/avalt/avsim) are available.\\
\texttt{outdegMixedPop} & behavior + another network + another network & nwsaom behavior(): only linear, quadratic, avalt, avsim (plus endow/creation of quadratic/avalt/avsim) are available.\\
\texttt{BFDeg} & behavior (non-directed network) + another network + another network & nwsaom behavior(): only linear, quadratic, avalt, avsim (plus endow/creation of quadratic/avalt/avsim) are available.\\
\bottomrule
\end{longtable}
\normalsize

\subsection*{Deliberately out of scope (80 rows)}
\noindent Continuous (SDE) dependent variables and their effects are not part of nwsaom (80 rows of the RSiena effects table).


